\documentclass[11pt]{article}
\usepackage{amsmath,amssymb,amsthm}
\newtheorem{theorem}{Theorem}\newtheorem{proposition}{Proposition}
\title{Native prefix paths, scalar probability and the output-order boundary}
\date{16 November 2026}
\begin{document}\maketitle
\paragraph{Audit and interface.} Written theorem with independent mathematical review; finite diagnostics are not formal verification. Fix an external input and a complete chosen controller for a locally finite normalized-row incremental arena, or its admitted feedback arena. All raw zero histories are retained. The primary observation is the coherent configuration/history prefix tower with its actual augmented tensor coefficients and full partial terminal bodies. The construction below adds a scalar Borel probability on its maximal histories; it neither forms an infinite coefficient tensor nor assigns an infinite polynomial grade. Controller tables, including off-policy and zero-history entries, remain separate structured data. Terminal means empty enabled-event set with its supplied terminal body; refusal, error and successful termination are not distinguished unless the arena supplies corresponding outcome tags.

\section{Finite raw prefix interface}
For $n\geq0$, let $H_n$ be the finite set of raw histories of $n$ selected scheduling steps, padding a terminal history by its unique absorbing continuation. Every selected nonterminal scope is nonempty and finite; its finite joint raw row has total mass one. All zero atoms remain in $H_n$. A zero-probability prefix still uses its declared normalized conditional table. Thus truncation $r_n:H_{n+1}\to H_n$ is surjective and has nonempty finite fibres, and the prefix masses satisfy
\[
 m_n(h)\geq0,\quad \sum_{h\in H_n}m_n(h)=1,
 \qquad m_n(h)=\sum_{r_n(k)=h}m_{n+1}(k).                 \tag{1}
\]
The finite-row assumption concerns each fixed input and horizon, not the size of the full weighted polynomial collection. It follows inductively from locally finite enabled sets and finite raw rows. If a separate finite normalized initial raw context is supplied, retain its actual raw atoms at $H_0$; otherwise $H_0$ is the singleton initial history. Let $T_n\subseteq H_n$ be the terminal subset. Terminal absorption gives $r_n^{-1}(T_n)\subseteq T_{n+1}$, not the reverse inclusion.

Define the all-raw path space
\[
 \Omega=\{(h_n)_{n\geq0}:r_n(h_{n+1})=h_n\}
 \subseteq\prod_{n\geq0}H_n,
 \qquad [h]=\{\omega:\omega_n=h\}.
\]
Give finite $H_n$ the discrete topology and $\Omega$ the induced product topology. It is compact and metrizable; cylinders form a countable clopen basis. Finite terminal histories are represented by their unique absorbing paths. The other paths are infinite active histories, including paths all of whose sufficiently long prefixes have zero mass.

\begin{theorem}[Probability without an infinite effect object]
There is a unique Borel probability $\mu$ on $\Omega$ such that $\mu([h])=m_n(h)$ for all $h\in H_n$. If $F_n$ is the event of first termination at step $n$, including initial termination $F_0$, and $D$ is the event of never terminating, then
\[
 \mu(D)=1-\sum_{n\geq0}\mu(F_n).                       \tag{2}
\]
\end{theorem}
\begin{proof}
Here is an explicit construction, avoiding any inference about coefficient limits. Partition $[0,1)$ into finitely many half-open intervals of lengths $m_0(h)$. Recursively partition the interval of each $h\in H_n$ into the finite child intervals of lengths $m_{n+1}(k)$ for $r_n(k)=h$. Equation (1) makes every partition exact; zero-length intervals are empty but their raw nodes remain in $\Omega$. At each level every point belongs to a unique nonempty interval, and its successive labels form a coherent path. This defines $f:[0,1)\to\Omega$. The inverse image of every cylinder is its assigned interval, hence $f$ is Borel measurable, because cylinders are a countable basis. Push forward Lebesgue probability. It has the stated cylinder values. No endpoint choice ambiguity arises with half-open intervals. A raw zero path need not be in the image of $f$; it is nevertheless present in the target space and has every appropriate zero cylinder measure.

Intersections of cylinders are either empty or cylinders, and cylinders generate the Borel sigma algebra. Two probabilities agreeing there agree on that sigma algebra by the standard elementary uniqueness argument: the sets on which the two measures agree form a Dynkin class containing the generating intersection-closed cylinder family. Finally $\{\omega:\omega_n\in T_n\}$ is increasing with $n$, its union is finite termination, and it is the disjoint union of the first-hit events. Countable additivity and complementation give (2).
\end{proof}

\section{Actual native maps and retained effects}
\begin{theorem}[Strict conditional-compatible naturality]
Suppose a native arena map, with a selected compatible controller pair, induces maps $f_n:H_n\to H'_n$ satisfying
\[
 r'_nf_{n+1}=f_nr_n,\qquad (f_n)_*m_n=m'_n,
 \qquad f_n^{-1}(T'_n)=T_n.                            \tag{3}
\]
Then $f_\infty(\omega)_n=f_n(\omega_n)$ is a continuous map of all-raw path spaces, $(f_\infty)_*\mu=\mu'$, and $f_\infty^{-1}(D')=D$. These constructions preserve identities and composition. An operational native isomorphism with transported controls gives a measure-preserving homeomorphism.
\end{theorem}
\begin{proof}
The first equation gives a coherent output path. A target cylinder pulls back to the finite union of source cylinders over its fibre, so the map is continuous. The second equation gives its cylinder masses; uniqueness of the preceding theorem gives the pushforward claim. The last equation gives equality of each finite termination event and hence divergence. Coordinatewise formulas establish functoriality and the inverse claim.
\end{proof}
For the intended native maps, (3) follows from the finite-prefix exact fibre equations, conditional-law compatibility at every raw past, and enabled-set bijections reflecting terminality. It is not asserted for arbitrary lax partial-body arrows or arbitrary initial-algebra folds without a chosen compatible controller. In particular, the labelled colored-feedback nesting isomorphism transports events, whole scopes and raw witnesses, hence satisfies (3) and gives the stated path isomorphism. Canonical backward controller lifting supplies another declared source of compatible pairs; no general forward controller descent is inferred.

The original actual augmented coefficient arrows at each $H_n$ remain part of the native prefix tower. Truncation deletes new event factors by their augmentations. The probability measure is only the scalar reduct of this structured object. A path determines a compatible sequence of its finite coefficient objects and deletion arrows, not a claimed object representing their infinite tensor. A terminal path retains its finite terminal tensor and full terminal grade. Nothing in (1)--(3) reconstructs coefficients, unused grade arrows, unreachable primitive rows, or off-policy controller tables from $\mu$.

\section{First-hit trace and output boundary}
In a guarded feedback arena the first-hit terminal raw histories of this tower are exactly the histories used by the admitted countable event-cost matching trace after its separately specified terminal/cost restriction. Partitioning finite terminal histories by first-hit step or by total number of paid primitive events does not change their disjoint union, their product weights, their actual finite event tensors, or their full terminal bodies. Thus (2) measures the complement of that terminating process's scalar mass. This is an equality after terminal aggregation, not an isomorphism of the unclosed step and cost clocks, and not a proof that general matching trace is least-fixed-point iteration. Under admitted well-founded silent initialization, finite silent closures change neither statement because they execute no primitive events and initial termination is included.

\begin{proposition}[Limit output need not be a stream]
Along a path, persistent finite emitted-port sets with order-preserving inclusions have a countable linearly ordered union. Its order need not be the natural-number order, or even a well-order. No sequential stream or infinite polynomial body follows from the finite prefix axioms.
\end{proposition}
\begin{proof}
The union is countable because it is a countable union of finite sets. Two ports eventually occur together, and their relative order is thereafter fixed, so the union order is well defined and total. Consider an arena with a single deterministic enabled event at every nonterminal step, each event used once and spawning the next, which inserts its new persistent output port before every previous port. Simultaneous-event diamond obligations are vacuous. The limit order is reverse natural-number order. More strongly, enumerate the rationals and at step $n$ emit the $n$th rational in its order relative to the previously emitted finite set. This is again a locally finite deterministic single-event arena, and the union order is the countable dense order of the rationals. Neither order is an ordinary sequential output list. Such an arena can have no terminal configurations, so its terminal-body obligation is vacuous; a separately supplied terminal branch can be added without changing this infinite path.
\end{proof}
A strict native path map with coherent port bijections induces an order isomorphism of the corresponding pathwise unions. That optional ordered-colimit observation remains separate from the full weighted finite terminal grade, from the scalar probability, and from any proposed infinite coefficient semantics. Applications needing conventional streams must add a condition such as append-only births and explicitly handle paths with finitely many outputs.
\end{document}
